\section{Introduction}
\label{sec:intro}

An agent that holds a relationship with a user over months accumulates a
transcript far larger than any context window, so every such system retrieves.
The dominant design is a dense retriever over conversation turns: embed each
turn, embed the question, return the nearest $k$. It works well enough that most
platforms ship it and report a single recall number.

A single recall number hides the failure. Questions that require \emph{one} fact
and questions that require \emph{several} fail in different ways, and averaging
them produces a figure that describes neither. We separate them, and the
separation is the paper.

On LoCoMo~\cite{locomo} --- ten multi-session conversations of 369--689 turns
with human-annotated evidence --- a dense retriever at $k{=}20$ retrieves every
required turn for \rSallTwenty{}\% of single-hop questions and \rMallTwenty{}\%
of multi-hop ones \meas{}. That gap looks like a retrieval-quality problem and
invites a retrieval-quality answer: a better embedding, a reranker, a larger $k$.

It is not a retrieval-quality problem. The same run reports a second metric that
settles it. Scoring \emph{any} rather than \emph{all} --- did the retriever find
at least one of the required turns --- gives \rSanyTwenty{}\% on single-hop and
\rManyTwenty{}\% on multi-hop, which are the same number. The retriever
understands multi-hop questions exactly as well as it understands single-hop
ones. It finds a relevant turn essentially always, and the rest of them
essentially never.

That is a precise diagnosis and it points at a known family of fixes. If one
search finds part of the answer, search again from what it found: pseudo-relevance
feedback~\cite{rocchio} in vector space, or the retrieve--reason--retrieve loop
of IRCoT~\cite{ircot}. This paper implements four such policies over a single
fixed index, measures them, and reports that \textbf{they do not work} --- and
then measures the corpus itself to explain why they cannot.

\subsection*{Contributions}

\begin{enumerate}[leftmargin=*]
\item A decomposition of dense-retrieval recall on long conversations into
  \emph{all}-evidence and \emph{any}-evidence, showing the two categories differ
  only in the former (\S\ref{sec:baseline}).
\item Four retrieval policies over one index, holding embedding, corpus and
  scoring fixed so that the policy is the only variable, with the negative result
  in full including the cases that got worse (\S\ref{sec:policies}).
\item A structural diagnostic of the evidence that explains the negative result
  and predicts it would survive a larger $k$, a better embedding, or a reranker
  (\S\ref{sec:diagnostic}).
\item A statement of what the measurements do and do not license anyone to build
  next (\S\ref{sec:implications}--\S\ref{sec:limits}).
\end{enumerate}
